% Generated by scripts/make_table.py from tables/evidence_matrix.csv. Do not edit by hand.
\begin{landscape}
\begingroup
\scriptsize
\singlespacing
\setlength{\tabcolsep}{3pt}
\setlength{\LTcapwidth}{\linewidth}
\renewcommand{\arraystretch}{1.15}
\begin{longtable}{>{\raggedright\arraybackslash}p{0.95in}>{\raggedright\arraybackslash}p{0.95in}>{\raggedright\arraybackslash}p{0.43in}>{\raggedright\arraybackslash}p{0.75in}>{\raggedright\arraybackslash}p{1.35in}>{\raggedright\arraybackslash}p{0.9in}>{\raggedright\arraybackslash}p{0.95in}>{\raggedright\arraybackslash}p{1.6in}>{\raggedright\arraybackslash}p{0.32in}}
\caption{Supplementary evidence matrix: experimental studies of rodent receivers cited in this review, in chronological order. Setting: lab; field; or semi (wild-caught animals tested in captivity or semi-natural conditions). nr, not reported in the sources consulted for this draft. Sec., manuscript section(s) where the study informs the argument (S2 features; S3 mechanisms; S4 action; S5 persistence; S6 integrated account). Limitations and verification status for each study are in tables/evidence\_matrix.csv.}\label{tab:evidence-supp} \\
\toprule
\textbf{Study} & \textbf{Species} & \textbf{Setting} & \textbf{Receiver experience} & \textbf{Stimuli} & \textbf{Controls} & \textbf{Endpoints} & \textbf{Key result} & \textbf{Sec.} \\
\midrule
\endfirsthead
\multicolumn{9}{@{}l}{\textit{Table~\ref{tab:evidence-supp}, continued}} \\
\toprule
\textbf{Study} & \textbf{Species} & \textbf{Setting} & \textbf{Receiver experience} & \textbf{Stimuli} & \textbf{Controls} & \textbf{Endpoints} & \textbf{Key result} & \textbf{Sec.} \\
\midrule
\endhead
\bottomrule
\endlastfoot
\citet{Davis1970} & Rat (Sprague-Dawley, male) & \mbox{lab} & Naive & Startle tones; Exp I: 2-s vs 16-s intervals (n=72); Exp II: regular vs variable at 8-s mean (n=144) & Schedule comparison & Startle amplitude; retest at 1 min and 24 h & 2-s spacing gave more decrement in training but 16-s spacing gave lower retest startle; regular timing gave more habituation than variable at matched mean & S5 \\
\rowcolor{black!5}\citet{DavisEtAl1982} & Rat & \mbox{lab} & Naive/\allowbreak{}startle & Acoustic startle pulses & Lesion/\allowbreak{}stimulation of startle pathway & Startle & Primary acoustic startle circuit includes PnC-related elements & S3 \\
\citet{EvansHammond1983a} & Rat (albino, male) & \mbox{lab} & Laboratory-reared & Recorded distress squeal at 80 or 100 dB; Exp 1: 10 trials 3 min apart (n=8 per group); Exp 2: retention after 0, 1, or 7 days (n=120); Exp 3: daily trials (n=32) & Synthetic mimic matched in burst timing & Orienting (lick suppression) & Natural squeal exceeded mimic at 100 but not 80 dB; all habituated rapidly; the 100-dB squeal recovered more after 1- and 7-day rests and habituated more slowly with 24-h spacing & S2\newline S5 \\
\rowcolor{black!5}\citet{EvansHammond1983b} & Rat & \mbox{lab} & Laboratory-reared & Natural distress squeal and mimic across contexts & Synthetic mimic & Orienting; generalization of habituation & Generalization of habituation across contexts depended on stimulus significance & S5 \\
\citet{LeDouxEtAl1984} & Rat & \mbox{lab} & Conditioned & Acoustic CS paired with shock & Lesion /\allowbreak{} pathway comparisons & Conditioned emotional responses & MGB/\allowbreak{}subcortical auditory projections contribute to conditioned responses to acoustic cues & S3 \\
\rowcolor{black!5}\citet{LeDouxEtAl1988} & Rat & \mbox{lab} & Conditioned & Acoustic CS + shock & Lesions of CeA projection targets & Arterial pressure; freezing & Distinct CeA targets dissociate conditioned autonomic vs freezing responses & S3\newline S6 \\
\citet{IwataLeDoux1988} & Rat & \mbox{lab} & Conditioned vs pseudoconditioned & Acoustic CS paired or unpaired with shock & Random, backward, shock-alone, naive groups; autonomic blockade & Arterial pressure; heart rate; freezing & Similar heart-rate rise in conditioned and control rats arose from sympathetic+parasympathetic coactivation vs sympathetic drive alone & S3 \\
\rowcolor{black!5}\citet{HendrieNeill1991} & Mouse (DBA/\allowbreak{}2, male) & \mbox{lab} & Laboratory-bred; predator-naive & 2-min calls of 10 species (predators: barn, tawny, little owl, kestrel, fox), 60-70 dB, after 3 min of human voice; n=10 per group & Human voice; non-predator calls (mallard, gull, oystercatcher, great tit, red deer) & Cage-, self-, and call-oriented behavior; defense; food intake & Barn and tawny owl calls increased self-oriented, call-oriented, and defensive behavior; kestrel and gull calls also increased defense; tawny owl call reduced feeding by about half (t(18)=2.67) & S2\newline S4\newline S3 \\
\citet{QuirkEtAl1995} & Rat & \mbox{lab} & Fear conditioned & Auditory CS & Pre/\allowbreak{}post conditioning & LA unit responses & Fear conditioning enhances short-latency LA auditory responses & S3 \\
\rowcolor{black!5}\citet{Mateo1996} & Belding's ground squirrel & \mbox{semi} & Juveniles; early auditory experience & Alarm calls (whistles, trills) & nr & Heart rate; behavior & Heart-rate discrimination of alarm calls preceded behavioral discrimination & S2\newline S4\newline S6 \\
\citet{LeeEtAl1996} & Rat & \mbox{lab} & Startle & Acoustic startle & Cochlear root /\allowbreak{} PnC interventions & Startle amplitude & Cochlear root neurons and PnC obligatory for primary startle pathway & S3 \\
\rowcolor{black!5}\citet{BlumsteinArmitage1997} & Yellow-bellied marmot & \mbox{field} & Free-living & Alarm calls varying in type, rate, and amplitude & nr & Vigilance and behavior & Calls tracked risk rather than predator type; receivers used call rate as risk information & S2\newline S3\newline S4\newline S5 \\
\citet{Hare1998} & Richardson's ground squirrel (juvenile females) & \mbox{field} & Free-living juveniles & Five calls about 3 min apart, each a repeated call of five chirps from a juvenile female, 84-91 dB; four of one exemplar, then a new exemplar from the same caller or a call from a different caller & Same caller, new exemplar & Vigilance duration; posture & Vigilance fell by the second call; a different caller restored it, a new exemplar from the same caller did not; neighbours' calls evoked longer vigilance than strangers' (trend, P=0.06) & S5 \\
\rowcolor{black!5}\citet{HendrieEtAl1998} & Voles, dormice, jerboas, spiny mice & \mbox{semi} & Wild-caught & Owl calls & nr & Behavior & Most species changed behavior; spiny mice responded less & S2\newline S4 \\
\citet{Shriner1998} & Yellow-bellied marmot; golden-mantled ground squirrel & \mbox{field} & Free-living & Conspecific and heterospecific alarm calls & Sympatric bird song & Vigilance & Both species distinguished alarms from song and responded similarly to conspecific and heterospecific alarms & S2\newline S4 \\
\rowcolor{black!5}\citet{EilamEtAl1999} & Gunther's vole; spiny mice & \mbox{semi} & Wild-caught & Owl calls & nr & Freezing or flight; corticosterone or cortisol & Voles froze or fled and raised corticosterone; spiny mice raised cortisol without behavioral change & S2\newline S4\newline S3\newline S6 \\
\citet{BaackSwitzer2000} & Eastern chipmunk & \mbox{field} & Free-living adults & Conspecific alarm calls & Control noise; silence & Vigilance; run to cover; re-emergence latency; seed intake & Alarms increased vigilance, shortened the run to cover, delayed re-emergence, and reduced food intake & S3\newline S4 \\
\rowcolor{black!5}\citet{HareAtkins2001} & Richardson's ground squirrel (juvenile females) & \mbox{field} & Free-living juveniles & Ten exemplars each of a reliable caller (paired with a stuffed badger except on playbacks 5 and 10) and an unreliable caller (no predator), 84-91 dB & Within-subject reliable vs unreliable caller & Postural vigilance; vigilance duration & Postural responses to the unreliable caller declined and differed from the reliable caller by playback 5 (P=.014) and 10 (P=.007); vigilance duration showed the same pattern only as a trend (P=.075) & S5\newline S6 \\
\citet{TsvetkovEtAl2002} & Rat & \mbox{lab} & Fear conditioned & Auditory fear conditioning & Slice physiology comparisons & LA synaptic plasticity & Fear conditioning occludes LTP-related enhancement at cortical input to LA & S3 \\
\rowcolor{black!5}\citet{MongeauEtAl2003} & Mouse & \mbox{lab} & With or without prior footshock & 17-20 kHz FM sweeps, 85 dB & Home vs unfamiliar cage; prior shock & Flight; freezing; neural activity & Cage familiarity and prior shock shifted responses between flight and freezing & S2\newline S3\newline S4 \\
\citet{BlumsteinDaniel2004} & Yellow-bellied marmot & \mbox{field} & Free-living & Habituation-recovery alarm playback; juvenile vs adult-female callers; familiar vs unfamiliar callers & Caller switch & Foraging suppression & Habituation to one caller recovered with a new caller; juvenile calls more evocative; familiarity null (Z=-0.30, p=.767) & S2\newline S5 \\
\rowcolor{black!5}\citet{BlumsteinEtAl2004} & Yellow-bellied marmot & \mbox{field} & Free-living; reliability manipulated & Alarms from reliable vs unreliable callers; degraded (distant) vs intact calls; one vs two callers & Baseline foraging; playback order & Foraging suppression; vigilance & Unreliable, distant-sounding, and two-caller playbacks were more evocative & S5 \\
\citet{SloanEtAl2005} & Richardson's ground squirrel & \mbox{field} & Free-living & Chirps, whistles, chucks & nr & Vigilance & Call types carried different information and altered vigilance differently & S2\newline S5 \\
\rowcolor{black!5}\citet{KiriazisSlobodchikoff2006} & Gunnison's prairie dog & \mbox{field} & Free-living (two colonies, 40-50 animals each) & 11-s alarm-call bouts originally elicited by humans, dogs, or coyotes; 45 dB at 50 m & Rubber-toy sound matched in rhythm, frequency range, and amplitude & First and second escape responses & Playbacks reproduced the escape responses to the corresponding live predators: burrow diving (humans), running to burrow and standing at rim (coyotes), alert posture in place (dogs) & S2\newline S3\newline S4 \\
\citet{SloanHare2006} & Richardson's ground squirrel & \mbox{field} & Free-living adults & Juvenile alarm calls with manipulated rate and bout length & Unmanipulated calls & Vigilance & Adults ignored rate and bout-length changes in juvenile calls & S5 \\
\rowcolor{black!5}\citet{EndresEtAl2007} & Rat (Sprague-Dawley, male; 59 total) & \mbox{lab} & Naive, then conditioned & Natural 22-kHz calls; 22-kHz sine tones; 22-kHz-range noise; other ultrasonic stimuli & Other frequency bands & Defensive behavior; freezing; conditioning and extinction & No single stimulus altered defense in naive rats (n=9); pooled 22-kHz range increased freezing; learned 22-kHz fear unusually persistent & S2\newline S4 \\
\citet{HerryEtAl2007} & Mouse (and human) & \mbox{lab} & Naive & 5-kHz, 20-ms, 85-dB pulses at 5 Hz (SD 140 ms) or 0.7 Hz (mean 1429 ms, SD 1368 ms) & Predictable vs jittered sequences & Basolateral amygdala units; c-Fos; avoidance & Jittered timing prevented rapid amygdala habituation, increased c-Fos, and promoted avoidance & S5 \\
\rowcolor{black!5}\citet{WohrSchwarting2007} & Rat (Wistar, male) & \mbox{lab} & Juvenile and adult cohorts & Natural 50-kHz calls; sine-wave approximations; 22-kHz calls & Background noise & Approach & Natural and synthetic 50-kHz calls both induced approach, especially in adults & S2 \\
\citet{SadanandaEtAl2008} & Rat (Wistar, male) & \mbox{lab} & Naive & 22-kHz calls; flat and FM 50-kHz calls & Background noise; silence & Approach; c-Fos & 50-kHz: approach and c-Fos in frontal cortex, accumbens, thalamus; 22-kHz: perirhinal cortex, amygdala, periaqueductal gray & S2\newline S3 \\
\rowcolor{black!5}\citet{SchmidtEtAl2008} & Eastern chipmunk & \mbox{field} & Free-living & Tufted titmouse mobbing calls; hawk calls & nr & Giving-up density & Foraging changed more to titmouse mobbing calls than to hawk calls & S2\newline S4 \\
\citet{SwanHare2008} & Richardson's ground squirrel & \mbox{field} & Free-living & Repeated calls; primary vs later syllables & nr & Vigilance & Primary syllables had especially strong alerting effects & S2\newline S5 \\
\rowcolor{black!5}\citet{BlumsteinRecapet2009} & Yellow-bellied marmot & \mbox{field} & Free-living & Alarm calls with added white-noise nonlinearities & Normal calls; silence-added calls & Foraging & Foraging fell more after calls with added nonlinearities & S2\newline S3\newline S4\newline S5 \\
\citet{KindermannEtAl2009} & Mouse (C57BL/\allowbreak{}6J), rat (Sprague-Dawley), pallid gerbil; all male & \mbox{lab} & Laboratory-reared, predator-naive & 3-min calls of owls and raptors, 70-80 dB peak, 0.3-9 kHz; Exp 2-3: calls as conditioned stimuli with footshock (10 pairings or 1) & Non-predator bird calls (chaffinch, blackbird, mallard, quail) & Freezing; exploration; vigilant rearing; conditioned freezing & No unconditioned defense to predator calls in any species (power to detect 12\% change in freezing); all calls conditionable after 10 pairings with no predator advantage; one-trial learning only for tawny owl, goshawk, and mallard & S2\newline S4 \\
\rowcolor{black!5}\citet{QuinnEtAl2009} & Rat & \mbox{lab} & Delay fear conditioned (2 or 16 trials) & Trained vs novel auditory CS & Post-training dorsal hippocampal lesion vs sham & Freezing & Dorsal hippocampal lesions reduced generalization to a novel tone after extended training & S3\newline S5\newline S6 \\
\citet{KimEtAl2010} & Rat (male pairs) & \mbox{lab} & Naive vs footshock-experienced receivers & 22-kHz calls from a fear-conditioned partner & Naive receivers; auditory thalamus lesion or inactivation & Freezing & Only experienced receivers froze; auditory thalamus required & S2\newline S4\newline S5\newline S6 \\
\rowcolor{black!5}\citet{Mateo2010} & Belding's ground squirrel & \mbox{semi} & nr & Trills (terrestrial predators) and whistles (aerial predators) & nr & Cortisol; heart rate & Trills raised cortisol; whistles were linked to slowed heart rate & S2\newline S4\newline S3 \\
\citet{AschemeierMaher2011} & Woodchuck; eastern chipmunk & \mbox{field} & Free-living & Conspecific and heterospecific alarm calls & nr & Vigilance & Woodchucks increased vigilance to both, more to conspecific alarms & S2\newline S4 \\
\rowcolor{black!5}\citet{SierraMercadoEtAl2011} & Rat & \mbox{lab} & Auditory fear conditioned & Conditioned tone during extinction & Muscimol inactivation of PL, IL, BLA, vHPC & Freezing; extinction memory & PL needed for fear expression, IL for extinction; BLA and vHPC for both & S3\newline S5 \\
\citet{HareWarkentin2012} & Richardson's ground squirrel & \mbox{field} & Free-living juveniles & Mother's vs other female's calls; colony vs non-colony callers & nr & Vigilance & No differential response to mother's or colony members' calls & S5 \\
\rowcolor{black!5}\citet{ParsanaEtAl2012a} & Rat (Sprague-Dawley) & \mbox{lab} & Naive & 22-kHz and 50-kHz calls & Frequency-matched tones & Freezing; 127 amygdala units & No freezing; matched tones reproduced most firing differences between call types & S2\newline S4\newline S3\newline S6 \\
\citet{ParsanaEtAl2012b} & Rat & \mbox{lab} & Footshock-experienced & 22-kHz calls & nr & Freezing & Proposed that rats learn from their own 22-kHz calls (autoconditioning) & S2\newline S4 \\
\rowcolor{black!5}\citet{GetschowEtAl2013} & Eastern gray squirrel & \mbox{field} & Free-living & Conspecific alarm; robin (sympatric), blackbird (allopatric, similar), honeyeater (allopatric, dissimilar) alarms & Avian control sounds including robin song & Vigilance; latency & Response varied with playback ($\chi^2$=19.68, p=.0014); robin and blackbird alarms like squirrel alarms; honeyeater less & S2\newline S4 \\
\citet{WilluhnEtAl2014} & Rat (male, food-restricted) & \mbox{lab} & Individually housed & Natural 50-kHz and 22-kHz calls, $\sim$69 dB, 15 s, 1-2 min intervals & Time- and amplitude-matched noise; background & Approach; accumbens dopamine (voltammetry) & 50-kHz calls increased approach and dopamine (n=12), both declining rapidly; 22-kHz and noise did not & S2\newline S3\newline S5\newline S6 \\
\rowcolor{black!5}\citet{JayneEtAl2015} & Eastern gray squirrel & \mbox{field} & Free-living & Predator calls during foraging & nr & Vigilance; foraging; distance to refuge & Increased vigilance, suppressed and delayed foraging; distance to refuge unchanged & S3\newline S4 \\
\citet{XiongEtAl2015} & Mouse & \mbox{lab} & Naive innate defense & Loud sound /\allowbreak{} sound-evoked flight assays & Corticocollicular perturbation & Flight & Auditory cortex to IC projections drive/\allowbreak{}amplify sound-evoked flight & S3 \\
\rowcolor{black!5}\citet{MarlinEtAl2015} & Mouse & \mbox{lab} & Virgin/\allowbreak{}maternal females & Pup calls & Oxytocin /\allowbreak{} cortical inactivation & Retrieval; cortical responses & Oxytocin in auditory cortex enables pup-call salience and retrieval plasticity & S3\newline S2 \\
\citet{GeisslerEtAl2016} & House mouse & \mbox{lab} & Experience with sound meaning & Sounds with context-specific meaning & nr & Auditory cortex activation & Cortical activation depended on the learned meaning of sounds & S3 \\
\rowcolor{black!5}\citet{TovoteEtAl2016} & Mouse & \mbox{lab} & Conditioned /\allowbreak{} optogenetic & Auditory fear conditioning; optogenetic activation & Cell-type specific manipulations & Freezing; flight interactions & CeA-vlPAG-medulla pathway mediates freezing; interacts with flight circuits & S3 \\
\citet{OudaEtAl2016} & Rat & \mbox{lab} & Playback receivers & Natural 22-kHz vs artificial call-like stimuli & Non-exposed & c-Fos in auditory and limbic areas & Natural 22-kHz engage amygdala/\allowbreak{}hippocampus more than artificial matched stimuli & S3 \\
\rowcolor{black!5}\citet{FadokEtAl2017} & Mouse & \mbox{lab} & Conditioned & Conditioned acoustic cues & Cell-type optogenetics/\allowbreak{}recordings & Freezing vs flight & Competing CeA populations support conditioned freezing vs flight & S3 \\
\citet{CalubEtAl2018} & Rat & \mbox{lab} & Footshock; laryngeal nerve transection vs sham & 22-kHz playback in a novel context & Sham surgery; vocalized vs not & Freezing & Preventing self-vocalization did not reduce acquired call-evoked freezing & S2\newline S4 \\
\rowcolor{black!5}\citet{ConcinaEtAl2018} & Rat (male) & \mbox{lab} & Auditory fear conditioned & Threat tone vs novel tone & Optogenetic inhibition of Te1 terminals in PL vs BLA; olfactory comparison & Freezing; LFP and multiunit synchrony & Inhibiting auditory cortex input to prelimbic cortex (not BLA) caused fear generalization; slow-gamma synchrony predicted discrimination & S3\newline S5\newline S6 \\
\citet{LillyEtAl2019} & Eastern gray squirrel & \mbox{field} & Free-living & Hawk call followed by bird chatter & Ambient noise & Vigilance & Vigilance relaxed faster with bird chatter & S2\newline S4\newline S5 \\
\rowcolor{black!5}\citet{GoodeEtAl2019} & Rat & \mbox{lab} & Conditioned & Forward vs backward /\allowbreak{} unpredictable auditory CS & BNST inactivation & Freezing & BNST required for freezing to poorly predictive threat signals & S3 \\
\citet{DemaestriEtAl2019} & Rat & \mbox{lab} & Playback receivers & 22-kHz vs 55-kHz USV playback & Silence & Anxiety-like behavior; c-Fos & Differential BLA/\allowbreak{}BNST c-Fos and behavior by call valence & S3 \\
\rowcolor{black!5}\citet{FreemanEtAl2020} & Richardson's ground squirrel (females, n=29) & \mbox{semi} & Laboratory chamber & Chirps or whistles from unfamiliar callers; 49 syllables over 2 min; 80 dB at 1 m & No-call control & c-Fos & Chirps reduced c-Fos in medial amygdala and superior colliculus; most regions unchanged & S3 \\
\citet{OlszynskiEtAl2020} & Rat (Wistar) & \mbox{lab} & Paired (n=30) vs single (n=31) housing & Natural 22-kHz and 50-kHz calls & Computer-generated tones & Vocalization; heart rate; locomotion & Natural calls usually produced larger effects, with similar directions; single housing enlarged several effects & S2\newline S3\newline S4\newline S6 \\
\rowcolor{black!5}\citet{SchiavoEtAl2020} & Mouse (female) & \mbox{lab} & Before vs after maternal experience & Pup-call temporal morphs (natural interval $\sim$175 ms) & Synthetic interval shifts & Pup retrieval; auditory cortex & Experience broadened and shifted behavioral and cortical tuning & S2\newline S4 \\
\citet{TasakaEtAl2020} & Mouse & \mbox{lab} & Maternal experience & Pup calls & nr & Temporal association cortex activity; retrieval & Temporal association cortex contributes to auditory-driven maternal plasticity & S2\newline S4\newline S3 \\
\rowcolor{black!5}\citet{WatermanMai2020} & Cape ground squirrel & \mbox{field} & Free-living & Crowned lapwing alarms; conspecific alarms & nr & Vigilance & As vigilant to lapwing alarms as to conspecific alarms & S2\newline S4 \\
\citet{PereiraEtAl2020} & Rat & \mbox{lab} & Shock-experienced receivers & Cessation of movement-evoked sound (silence) & Continuous movement sound; optogenetic inactivation & Freezing & MGD, ventral auditory cortex, and LA necessary for freezing to silence cue & S3 \\
\rowcolor{black!5}\citet{HernandezEtAl2021} & Abrothrix; Oligoryzomys; Rattus rattus & \mbox{field} & Free-living & Calls of two owl species & No owl call & Feeding; locomotion; vigilance (camera traps) & Predator- and species-specific changes; moonlight and cover also mattered & S2\newline S4\newline S3 \\
\citet{VielleEtAl2021} & Rat & \mbox{lab} & Cagemates and strangers & Self-administered playback of cagemate vs stranger 50-kHz calls & Subthalamic lesion vs sham & Playback choice; cocaine intake & Preferred stranger calls; only stranger calls reduced cocaine intake; lesions abolished both & S2\newline S5 \\
\rowcolor{black!5}\citet{LiEtAl2021} & Mouse & \mbox{lab} & Naive & Auditory looming (crescendo noise) & Receding controls; pathway-specific perturbations & Freezing then flight & Corticostriatal (TS) vs corticocollicular (SC) pathways dissociate freezing vs flight & S3 \\
\citet{BerzEtAl2022} & Rat & \mbox{lab} & Playback followed by social interaction & 50-kHz calls & nr & Approach; brain activation & Approach habituated despite contingent social interaction & S5 \\
\rowcolor{black!5}\citet{GoedhoopEtAl2022} & Rat (Long-Evans, male) & \mbox{lab} & Conditioned & White noise 70-96 dB (typically 90 dB, 6 s) predicted by a 5-s cue & Intensity; probability & Accumbens core dopamine (voltammetry, n=21); avoidance; foraging; operant choice & Dopamine fell to noise and cue; did not scale with intensity (p=.88); unexpected noise slowed recovery & S3\newline S5 \\
\citet{GranataEtAl2024} & Rat & \mbox{lab} & Awake restrained fMRI & 22-kHz /\allowbreak{} 55-kHz USV; synthetic 22-kHz tone & Silence baseline & BOLD (BNST, BLA) & Valence-related limbic BOLD during USV playback; tone weaker & S3 \\
\rowcolor{black!5}\citet{WarrenEtAl2025} & Prairie vole (female) & \mbox{lab} & Pair-bonded & Partner vs stranger male calls, including calls directed at another female & nr & Preference behavior & Discriminated partner from stranger calls & S2\newline S5 \\
\citet{DuanEtAl2025} & Mouse & \mbox{lab} & Auditory threat conditioning & Conditioned auditory CS & LC/\allowbreak{}NE and dmPFC manipulations & Freezing; unit/\allowbreak{}photometry & LC-NE needed to convert retrieved auditory fear memory into freezing & S3 \\
\rowcolor{black!5}\citet{FuruyamaEtAl2026} & Mouse (C57BL/\allowbreak{}6 and DAT-Cre, both sexes) & \mbox{lab} & Head-fixed & Noise bursts 45-105 dB, 10 s; ramped and damped noise & Intensity series & Accumbens and tail-striatum dopamine (photometry); striatal neuron calcium; pupil; place avoidance at 85 dB & Accumbens dopamine fell more with louder noise (9 decrease-type mice); 6 others also showed an onset increase on the first burst; tail-striatum dopamine rose & S3 \\
\end{longtable}
\endgroup
\end{landscape}
